% Problem dossier: VI.06.P1
\subsection*{Problem VI.06.P1 --- Ideal containment as geometric incidence}
\label{prob:vi-06-p1}

\paragraph{Problem.}
Let \(A\) be a ring, \(I\subseteq A\) an ideal, and \(\mathfrak p\subseteq A\) a prime
ideal.
\begin{enumerate}
\item Prove that if \(I=(f_1,\dots,f_r)\), then
\[
I\subseteq\mathfrak p
\quad\Longleftrightarrow\quad
f_1,\dots,f_r\in\mathfrak p.
\]
\item Compute the prime ideals containing \((6)\subset\mathbb Z\).
\item In \(k[x,y]\), show that both \((x)\) and \((y)\) contain \((xy)\), and explain why every prime containing \((xy)\) contains one of them.
\item Over algebraically closed \(k\), interpret \(I\subseteq\mathfrak p\) as incidence of the irreducible locus \(V(\mathfrak p)\) with the closed condition \(V(I)\).
\end{enumerate}

\ifdefined\FullProblemDossiers
\paragraph{Why this problem matters.}
The definition of a closed set in a prime spectrum is built from ideal containment.  Before
introducing that topology, this problem makes the underlying incidence relation completely
concrete.

\paragraph{Useful ideas before starting.}
Generated ideals, divisibility in \(\mathbb Z\), the prime-product property, and order
reversal under \(V\).

\paragraph{Guided hints.}
For (1), an ideal contains the generators if and only if it contains every linear
combination of them.  For (2), use divisibility.  For (3), apply primality to \(xy\).  For
(4), use \cref{prop:vi06-incidence}.

\paragraph{Solution.}
For (1), if \(I\subseteq\mathfrak p\), all generators lie in \(\mathfrak p\).  Conversely,
if every generator lies in \(\mathfrak p\), then every sum \(\sum a_i f_i\) lies there
because \(\mathfrak p\) is an ideal.

For (2), the primes containing \((6)\) are \((2)\) and \((3)\): a prime ideal \((p)\)
contains \((6)\) exactly when \(p\mid6\).

For (3), \(xy\) lies in both \((x)\) and \((y)\).  If \(\mathfrak p\) contains \((xy)\),
then \(xy\in\mathfrak p\), so primality gives \(x\in\mathfrak p\) or
\(y\in\mathfrak p\).  Thus \((x)\subseteq\mathfrak p\) or \((y)\subseteq\mathfrak p\).

For (4), over algebraically closed \(k\),
\[
I\subseteq\mathfrak p
\quad\Longleftrightarrow\quad
V(\mathfrak p)\subseteq V(I).
\]
So the prime point lies on the closed condition defined by \(I\) precisely when the entire
irreducible locus represented by that prime satisfies those equations.

\paragraph{What happened mathematically.}
Membership of a future spectrum point in a closed set will be nothing more than containment
of ideals.  Geometry is therefore encoded by the partial order of ideals.

\paragraph{Extensions.}
Repeat the arithmetic computation for \((n)\subset\mathbb Z\), and describe the answer in
terms of the distinct prime divisors of \(n\).

\paragraph{Provenance.}
Recovered and reorganized from the legacy section ``Can a prime ideal contain another
ideal?'' in \texttt{theory-of-commutative-algebra-13.tex}.
\fi
